\section{What Comes Next}

The platform covers a lot of ground already, but there are clear directions where it needs to grow.

\subsection{Connecting Infrastructure to Observing Conditions}

Right now the platform is good at telling you what the infrastructure is doing. What it does not
do well yet is connect that to what the observatory is doing: which part of the observing sequence
is running, what the weather is, what state the camera and telescope are in, whether a calibration
is happening or a filter just changed. Some of that context should come in as labels. Some of it
may work better as events on a shared timeline. Either way, an operator trying to understand why
something degraded at a specific moment needs both views together, not separately.

\subsection{Tracing}

Metrics and logs cover most situations well, but there are cases---particularly in pipelines that
cross multiple services, queues, and storage systems---where it is hard to find exactly where a
request started failing. Distributed tracing can help with that \cite{opentelemetry-traces}. It is
not something to add everywhere, because it comes with instrumentation overhead and more data to
store, but in the right places it could significantly reduce the time spent finding the first
failing boundary.

\subsection{AI-Assisted Analysis}

One recent step that is already showing promise is enabling API access to Loki for developers
across the observatory. What started as a practical decision has opened up uses that go well beyond
what the dashboards were designed for. Teams are already building tools that use AI to analyze log
patterns, surface anomalies, and find correlations that would be extremely difficult to catch by
eye. The early results have been impressive. This is still the beginning, and it is hard to predict
exactly where it goes, but the direction is clear: when the data is accessible and the labels are
good, people find ways to ask questions the platform was never explicitly built to answer.